% concatenated from main.tex, shortcuts.tex
% Original pdfTeX output-mode primitive omitted for XeTeX.
\documentclass[10pt,a4paper]{amsart}

\setlength{\textwidth}{\paperwidth}
\addtolength{\textwidth}{-1in}
\setlength{\textheight}{\paperheight}
\addtolength{\textheight}{-1in}
\calclayout

\newcommand\commentm[1]{{\color{magenta} \sf [#1]}}
\newcommand{\JC}[1]{{\color{magenta}\sf $\clubsuit\clubsuit\clubsuit$ Jerson: #1}} 
\newcommand{\FP}[1]{{\color{blue}\sf $\clubsuit\clubsuit\clubsuit$ Fabien: #1}}
\newcommand{\RP}[1]{{\color{teal}\sf $\clubsuit\clubsuit\clubsuit$ Riccardo: #1}}

\usepackage[utf8]{inputenc}
    \usepackage[T1]{fontenc}
    \usepackage[sc]{mathpazo}
    \linespread{1.05}
    
    \usepackage{stackengine,amsmath,amsthm,amssymb,mathtools,soul,wrapfig}
    \usepackage{mathdots}
    \usepackage{graphicx}
    \usepackage{tikz,tikz-cd}
            \tikzset{
                subseteq/.style={
                draw=none,
                edge node={node [sloped, allow upside down, auto=false]{$\subseteq$}}},
                Subseteq/.style={
                draw=none,
                every to/.append style={
                edge node={node [sloped, allow upside down, auto=false]{$\subseteq$}}}
                }
            }
            \tikzset{
                equal/.style={
                draw=none,
                edge node={node [sloped, allow upside down, auto=false]{$=$}}},
                Equal/.style={
                draw=none,
                every to/.append style={
                edge node={node [sloped, allow upside down, auto=false]{$=$}}}
                }
            }
            \tikzset{
            rotated_label/.style={anchor=south, rotate=90, inner sep=.5mm}
            }
    \usetikzlibrary{decorations}
    \usepackage{verbatim}
    \usepackage{booktabs}
    \usepackage{array}
    \usepackage{tabularx}
    \usepackage{booktabs}
    \usepackage{siunitx}
    \usepackage{rotating}
    \usepackage{adjustbox}
    \usepackage{longtable}
    \usepackage{caption} 
    \usepackage{stmaryrd}
    \usepackage{nicefrac}
    \usepackage{enumitem}
    \usepackage{xfrac}
    \usepackage[textwidth=0.8in,textsize=tiny]{todonotes}
        \setlength{\marginparwidth}{0.8in}
    
\usepackage[backref=true,giveninits=true,doi=false,url=false,isbn=false,backend=biber]{biblatex}
    \newbibmacro{string+doiurlisbn}[1]{%
      \iffieldundef{doi}{%
        \iffieldundef{url}{%
          \iffieldundef{isbn}{%
            \iffieldundef{issn}{%
              #1%
            }{%
              \href{http://books.google.com/books?vid=ISSN\thefield{issn}}{#1}%
            }%
          }{%
            \href{http://books.google.com/books?vid=ISBN\thefield{isbn}}{#1}%
          }%
        }{%
          \href{\thefield{url}}{#1}%
        }%
      }{%
        \href{http://dx.doi.org/\thefield{doi}}{#1}%
      }%
    }
    \DeclareFieldFormat{title}{\usebibmacro{string+doiurlisbn}{\mkbibemph{#1}}}
    \def\CalO{{\mathcal{O}}}
    \def\CalQ{{\mathcal{Q}}}
    \def\CalH{{\mathcal{H}}}
    \def\CalS{{\mathcal{S}}}
    \DeclareMathOperator*{\Res}{Res}
    \DeclareFieldFormat[article,incollection,thesis,misc,inproceedings,online,inbook]{title}{\usebibmacro{string+doiurlisbn}{\mkbibquote{#1}}}
    \DeclareFieldFormat{year}{\usebibmacro{year-parenthesis}{\mkbibemph{#1}}}
\renewbibmacro*{date}{%
\iffieldundef{year}
{}
{\ifentrytype{misc}{\printtext[parens]{\printdate}}{\printdate}}}%
    \addbibresource{mybib.bib}

\usepackage{hyperref}
    \hypersetup{colorlinks=true,allcolors=blue}
    \usepackage[capitalize,nameinlink,noabbrev]{cleveref}

\usepackage{amsthm}
    \theoremstyle{plain}
        \newtheorem{theorem}{Theorem}[section]
        \newtheorem{lemma}[theorem]{Lemma}
        \newtheorem{proposition}[theorem]{Proposition}
        \newtheorem{corollary}[theorem]{Corollary}

    \theoremstyle{definition}
        \newtheorem{definition}[theorem]{Definition}

    \theoremstyle{remark}
        \newtheorem{example}[theorem]{Example}
        \newtheorem{notation}[theorem]{Notation}
        \newtheorem{remark}[theorem]{Remark}
        \newtheorem{question}[theorem]{Question}


% Circled numbers
\newcommand*\circled[1]{\tikz[baseline=(char.base)]{
            \node[shape=circle,draw,inner sep=2pt] (char) {#1};}}
        
 \DeclareMathOperator{\GL}{GL}
\DeclareMathOperator{\SL}{SL}
\DeclareMathOperator{\PSL}{PSL}
\DeclareMathOperator{\Gal}{Gal}
\DeclareMathOperator{\lcm}{lcm}
\DeclareMathOperator{\disc}{disc}
\DeclareMathOperator{\car}{char}
\DeclareMathOperator{\End}{End}
\DeclareMathOperator{\Aut}{Aut}
\DeclareMathOperator{\red}{red}
\DeclareMathOperator{\id}{id}
\DeclareMathOperator{\ord}{ord}
\DeclareMathOperator{\Split}{Split}
\DeclareMathOperator{\comp}{Compositum}
\DeclareMathOperator{\sgn}{sgn}
\DeclareMathOperator{\Frob}{Frob}

\newcommand{\A}{\mathcal{A}}
\newcommand{\B}{\mathcal{B}}
\newcommand{\F}{\mathbb{F}}
\newcommand{\M}{\mathfrak{M}}
\newcommand{\K}{\mathfrak{K}}
\newcommand{\Sing}{\mathfrak{S}}
\newcommand{\N}{\mathbb{N}}
\newcommand{\Z}{\mathbb{Z}}
\newcommand{\Q}{\mathbb{Q}}
\newcommand{\R}{\mathbb{R}}
\newcommand{\C}{\mathbb{C}}
\newcommand{\U}{\mathcal{U}}
\newcommand{\PP}{\mathcal{P}}
\newcommand{\OO}{\mathcal{O}}
\newcommand{\PPP}{\mathbb{P}}
\newcommand{\ii}{\mathrm{i}}
\newcommand{\eps}{\varepsilon}
\newcommand{\li}{\operatorname{li}}
\newcommand{\mcm}{\operatorname{lcm}}
\newcommand{\card}{\operatorname{card}}
\newcommand{\diag}{\operatorname{diag}}
\newcommand{\res}{\operatorname{res}}
\newcommand{\shiftLone}{\operatorname{\tt shiftL1}}
\newcommand{\shiftLeight}{\operatorname{\tt shiftL8}}
\newcommand{\inv}{\operatorname{\tt inv}}
\newcommand{\dx}{\mathrm{d}}
\newcommand{\Zn}[1]{\Z_{#1}}
\newcommand{\Zs}[1]{\Zn{#1}^*}
\newcommand{\eqdef}{\mathchoice{{\,\overset{\mathrm{def}}=}\,}{:=}{:=}{:=}}
\newcommand{\Odip}[2]{\mathcal{O}_{#1}\!\left(#2\right)\mathchoice{\!}{}{}{}}
\newcommand{\Odi}[1]{\Odip{}{#1}}
\newcommand{\vuoto}{\textup{\textrm{\O}}}
\newcommand{\Req}{\mathcal{R}}
\newcommand{\REQ}{\mathrel{\Req}}
\newcommand{\cstylemod}{\,\mathbin{\%}\,}
\newcommand{\cstyleinc}{\,\mathbin{+\!\!=}\,}
\newcommand{\cstyledec}{\,\mathbin{-\!\!=}\,}
\newcommand{\cstyleplusplus}{\,\mathbin{+\!\!+}\,}
\newcommand{\legendre}[2]{\mathchoice{\genfrac{(}{)}{}0{#1}{#2}}
  {\bigl(#1\mid#2\bigr)}{(#1\mid#2)}{(#1\mid#2)}} 
\newcommand{\Ii}{\mathcal{I}}
\newcommand{\spazio}{\char`\ }
\newcommand{\ttspc}{\texttt{\spazio}}
\newcommand{\Eulerphi}{\varphi}
\newcommand{\Dgotic}{\mathfrak{D}}
\newcommand{\Pgotic}{\mathfrak{P}}
\newcommand{\pgotic}{\mathfrak{p}}
\newcommand{\qgotic}{\mathfrak{q}}
\newcommand{\Kbar}{\overline{K}}
\newcommand{\kbar}{\overline{k}}
\newcommand{\kp}{k_{\mathfrak{p}}}
\newcommand{\fcap}{\hat{f}}
\newcommand{\gcap}{\hat{g}}
\newcommand{\Ogotic}{\mathcal{O}}
\newcommand{\rgotic}{\mathfrak{r}}
\newcommand{\m}{\mathfrak{m}}
\newcommand{\fgotic}{\mathfrak{f}}
\newcommand{\fgoticO}{\mathfrak{f}_\mathcal{O}}
\newcommand{\bgotic}{\mathfrak{b}}
\newcommand{\OgoticK}{\mathfrak{O}_K}
\newcommand{\den}{\text{den}}
\newcommand{\Qbar}{\overline{\Q}}
\newcommand{\trdeg}{\text{trdeg}}
\newcommand{\rad}{\text{rad}}
\newcommand{\xitilde}{\tilde{\xi}}
\newcommand{\xvect}{\textbf{x}}
\newcommand{\LF}{\mathcal{L}(\textit{F})}


\newcommand{\norm}[1]{\lVert#1\rVert}
%\renewcommand{\v}[1]{\ensuremath{\mathbf{#1}}}
\newcommand{\dist}{\operatorname{dist}}
\newcommand{\rot}{\operatorname{rot}}
\renewcommand{\mod}{\operatorname{mod}\ }
\newcommand{\rk}{\operatorname{rank}}

\newcommand\Tstrut{\rule{0pt}{2.9ex}}         % = `top' strut
\newcommand\Bstrut{\rule[-1.65ex]{0pt}{0pt}}   % = `bottom' strut
\newcommand{\llrrparen}[1]{% \llrrparen{..}
  \left(\mkern-6mu\left(#1\right)\mkern-6mu\right)}
\newcolumntype{C}[1]{>{\centering\arraybackslash}p{#1}}

% Function restriction
\newcommand\restr[2]{{ \left.\kern-\nulldelimiterspace #1 \vphantom{\big|} \right|_{#2} }}

\newcommand\altxrightarrow[2][0pt]{\mathrel{\ensurestackMath{\stackengine%
  {\dimexpr#1-7.5pt}{\xrightarrow{\phantom{#2}}}{\scriptstyle\!#2\,}%
  {O}{c}{F}{F}{S}}}}

%\renewcommand{\baselinestretch}{1.2}
%\overfullrule=1mm
\raggedbottom

\makeatletter
\@namedef{subjclassname@2020}{\textup{2020} Mathematics Subject Classification}
\makeatother

\hyphenation{Uni-ver-sit-ets-park-en}
\AtBeginDvi{\special{pdf:mapfile upl.map}\special{pdf:mapfile pazo.map}\special{pdf:mapfile fpls.map}}
